\section{Related Work}
\label{sec:related_work}

Our research lies at the intersection of weight-space model merging, post-training network quantization, and test-time gradient-free and gradient-based adaptation.

\subsection{Weight-Space Model Merging}
Model merging has emerged as a key zero-shot method to create unified multi-task networks without the prohibitive computational cost of joint multi-task training~\cite{ilharco2022editing, wortsman2022model}. Early approaches averaged weights directly, forming \emph{Model Soups} for improved generalization on a single task~\cite{wortsman2022model}, or utilizing \emph{Task Arithmetic} to add task-specific parameter updates (task vectors) to a shared pre-trained base model~\cite{ilharco2022editing}. Subsequent works focused on resolving weight-space interference through coordinates masking or sign-alignment (e.g., TIES-Merging~\cite{yadav2023resolving} and DARE~\cite{yu2023language}), or analyzing linear mode connectivity modulo permutation symmetries (Git Re-Basin~\cite{ainsworth2022git, entezari2021role}).

To move beyond static uniform merging coefficients, recent methods have introduced test-time adaptive coefficient optimization. \emph{AdaMerging}~\cite{yang2024adamerging} optimizes layer-wise coefficients on unlabeled data streams via entropy minimization, while \emph{SyMerge}~\cite{jung2025symerge} optimizes coefficients to induce positive task synergy. However, these methods operate exclusively in full-precision (FP16 or FP32) weight spaces. When practitioners attempt to deploy these merged models to edge hardware, subsequent quantization degrades performance. Q-Merge is the first framework to formulate and solve model merging directly under the quantization operator.

\subsection{Post-Training Quantization (PTQ)}
Quantization is the industry-standard method for compressing deep learning models to reduce their memory footprint and accelerate inference on edge devices~\cite{jacob2018quantization, gholami2021survey}. Unlike Quantization-Aware Training (QAT), which requires expensive full-training loops with labeled data, Post-Training Quantization (PTQ) compresses pre-trained models using either no data or a small calibration subset, making it highly practical for deployment~\cite{nagel2020adaround}. 

While 8-bit PTQ (INT8) is highly robust and widely supported, aggressive low-bit configurations (e.g., INT4) often suffer severe accuracy degradation due to massive quantization noise. Advanced PTQ techniques attempt to mitigate this noise by learning adaptive rounding offsets (AdaRound~\cite{nagel2020adaround}), scaling weight and activation ranges (SmoothQuant~\cite{xiao2023smoothquant}, AWQ~\cite{lin2023awq}), or utilizing second-order Taylor expansions (GPTQ~\cite{frantar2022gptq}). Despite these advances, prior PTQ research is restricted to single-task models. Q-Merge addresses the unique challenge of PTQ in multi-task merged networks, where the quantization operator acts on a weight space that is dynamically parameterized by task-merging coefficients.

\subsection{Straight-Through Estimator and Optimization under Discretization}
Optimizing model parameters under discrete rounding constraints is notoriously difficult because the rounding function is a step function with zero gradient almost everywhere and infinite gradient at step boundaries, which completely breaks standard backpropagation. The Straight-Through Estimator (STE), introduced by Bengio et al.~\cite{bengio2013estimating}, addresses this by approximating the gradient of the round function as an identity mapping during the backward pass ($\frac{\partial \text{round}(x)}{\partial x} \approx 1$). STE has been highly successful in binary neural networks and QAT~\cite{jacob2018quantization}. 

An alternative to gradient-based estimation in non-differentiable landscapes is derivative-free black-box optimization, such as Evolutionary Strategies (1+1 ES) or random mutation sweeps~\cite{akiba2024evolutionary, yang2024adamerging}. While zero-order methods bypass non-differentiability, they often exhibit high variance and slow convergence in high-dimensional spaces. In this work, we present a rigorous comparison of STE-based first-order gradient descent and zero-order 1+1 ES stochastically optimizing merging coefficients, exposing the superior stability and efficiency of gradient-guided coordinates alignment in quantized weight spaces.

\subsection{Low-Bit Model Merging and Task Vector Compression}
Very recently, several researchers have begun exploring the intersection of low-bit quantization and model merging to enable memory-efficient multi-task deployment on edge hardware. For instance, Task Vector Quantization (TVQ)~\cite{tvq2025} and Residual Task Vector Quantization (RTVQ) address the storage burden of carrying multiple fine-tuned models on-device by quantizing the task vectors (parameter updates) directly to ultra-low bitwidths (e.g., 2-bit), leaving the pre-trained base model in full precision. Similarly, 1bit-Merging~\cite{onebitmerging2025} employs dynamic quantization scales across key layers to compress task vectors while leveraging a trained router to manage task-specific weights during inference.

Furthermore, Merge-Friendly Post-Training Quantization via Hessian and Distant Regularizing Quantization (HDRQ)~\cite{hdrq2025} tackles multi-domain adaptation by flattening the loss landscape during quantization. This ensures that when multiple quantized experts are merged, they do not collide with steep error barriers that degrade performance. Additionally, E-PMQ~\cite{epmq2026} uses source expert weights as anchors during post-merge quantization to protect against structural merging deviation.

While these approaches offer elegant solutions for static weight-space compression or expert-guided rounding, they either require full-precision base checkpoints during inference (like TVQ), pre-merging optimization of separate models, or do not support test-time adaptive coefficient blending under joint quantization noise. Q-Merge conceptually differs by and complements these works: it optimizes the blending coefficients $\Lambda$ directly under the joint quantization operator at test-time on a small, unlabeled stream, and can be sequentially combined with methods like HDRQ or TVQ to find the optimal coordinate-alignment before task-vector compression is executed.